\scanpage{69}
Cas des paires $\{x,y\}$ de pts dist. non liés. La donnée équivaut à une \uline{bibase} $B=\{b,b'\}$ avec un elt $z$ de $B/\sigma_B=I$. La catégorie des $(C,\{x,y\})$ équivaut à celle des couples $(\tau,J)$ de deux ens avec $\tau$ de cardinal $2$, \inserted{$(\tau=\{b,b'\})$} $J$ de cardinal $5$ ($J=B/\sigma_B-\{z\}$).
\margin{$\unsure{\widehat{\mathrm{Ens}}_2|\mathrm{Ens}_2}\times\widehat{\mathrm{Ens}}_5$}
% slanted
\noindent\rule{2.3cm}{.4pt}

$\uline{\mathrm{Pl}}(L_3,C)$. Le type d'une partie \inserted{épinglée} trois-libre \inserted{$(x,y,z)$} de $C$ est $(\unsure{\delta_{12}|J_{12}|\xi_{12}},\delta_{23},\delta_{31})$. L'ens des pts de $C-\{\{\delta_{12},\delta_{23},\delta_{31}\}$ qui ne sont \deleted{pas} liés \deleted{aux} \inserted{d'aucun} des $x,y,z$ est $\{\underbrace{\delta_4,\delta_5,\delta_6},\underbrace{\delta'_{\unsure{45|4}},\delta'_{\unsure{56|5}},\delta'_{\unsure{65|6}}}\}$ qui est de

cardinal $6$; le groupe des autom. de $(C;\delta_{12},\delta_{23},\delta_{31})$ (contenant le groupe des autom de $(C,\{\{\delta_i\}\},\{\delta'_i\})$ qui respecte la partition \inserted{de $b\unsure{=|\simeq}$} $[1,6]$ en $[1,3]$ et $[4,6]$) est transitif sur l'ens précédent.

\uline{Prop} Pour tt couple \inserted{$(B,K)$} formé d'une bibase $B=\{b,b'\}$ \inserted{de $C$} et d'une partie $K$ de $B/\sigma_B$ telles que $\mathrm{Card}\,K=3$, soit $J(B,K)$ l'ens partie de $C$ formée les $\delta^B_{ij}$ $(i,j\in K)$. On trouve ainsi une bijection de l'ens \deleted{\unsure{d|l|$\mathcal B$}} de cardinal $72.\underbrace{\binom{6}{3}}_{20}\frac{1}{2}=720)$ des couples $(B,K)$ avec l'ens des parties $3$-libres de $C$.

La cat. des couples $(C,J)$ d'un graphe cubique et d'une partie $3$-libre de $C$, équivaut donc à celle des triplets $(\tau,K,K')$ d'ens \margin{$\widehat{\mathrm{Ens}}_2\times\widehat{\mathrm{Ens}}_3\unsure{\times|-}\widehat{\mathrm{Ens}}_3$} de cardinal $2,3,3$ -- en associant à un triple le complexe $C=\unsure{\mathbb E|\mathscr E|C}(\tau,K\amalg K')=(\tau\times(K\amalg K'))\amalg\unsure{\mathfrak P_2|\mathscr P_2}(K\amalg K')$, et l'ens $J$ des $\delta^B_{ij}$ $(i,j\in K)$, où $B$ est la bibase can. de $C$. Le groupe
% slanted
% archive sheet number: 68
